\documentclass{article}
\usepackage{amsmath, amssymb, amsthm}
\usepackage{hyperref}
\usepackage{geometry}
\geometry{margin=1in}

\title{Erd\H{o}s Problem \#137}
\date{Last updated: 2025-08-31}
\author{Source: erdosproblems.com}

\begin{document}
\maketitle

\noindent\textbf{Status:} OPEN \quad \textbf{No prize}

\noindent\textbf{Tags:} number theory

\medskip

\noindent\textbf{Problem Statement:}

We say that $N$ is powerful if whenever $p\mid N$ we also have $p^2\mid N$. Let $k\geq 3$. Can the product of any $k$ consecutive positive integers ever be powerful?




Conjectured by Erd\H{o}s and Selfridge. There are infinitely many $n$ such that $n(n+1)$ is powerful (see [364] ). Erd\H{o}s and Selfridge \cite{ErSe75} proved that the product of $k\geq 3$ consecutive positive integers can never be a perfect power. Erd\H{o}s remarked that this 'seems hopeless at present'. In \cite{Er82c} he further conjectures that, if $k$ is fixed and $n$ is sufficiently large, then, for all $m$, there must be at least $k$ distinct primes $p$ such that\[p\mid m(m+1)\cdots (m+n)\]and yet $p^2$ does not divide the right-hand side. See also [364] .




References

[Er82c] Erd\H{o}s, P., Miscellaneous problems in number theory . Congr. Numer. (1982), 25-45.

[ErSe75] Erd\H{o}s, P. and Selfridge, J. L., The product of consecutive integers is never a power . Illinois J. Math. (1975), 292-301.


\bigskip
\noindent\small{Source: \url{https://www.erdosproblems.com/137}}
\end{document}
